%%% FIELD proof-prop-one-sided-clinical-followup-rule
The two source-record propositions supply the factual premises about Liberali et al.  Specifically, \(\texttt{prop:liberali-main-source-record}\) and \(\texttt{prop:liberali-supplement-source-record}\) state that the published analysis is a retrospective resampling experiment: conditional on the arm selected by the simulated real-time rule, a patient is drawn with replacement from the corresponding original-trial arm pool; the experiment is repeated for \(200\) replicates; and the reported lower and upper uncertainty limits are the empirical \(2.5\%\) and \(97.5\%\) replicate quantiles.  The same records state that GUSTO-1 contributes three analyzed arms and a death-or-survival endpoint at \(30\) days, whereas EUROPA contributes two arms and a composite cardiovascular endpoint, with a reported \(4.2\)-year mean time to the last observation.  They also state that the baseline simulated rule makes daily Gittins-index assignments and that equal randomization in the eta-variant applies only while an arm has not reached its fixed minimum count.

These reported percentile intervals are different inferential maps from both intervals defined in this note.  Indeed, by \(\texttt{def:ordinary-wald}\), the ordinary-Wald map applied to one logged experiment is
\[
 I_{T,\Delta}^{\mathrm W}
 =\bigl[\widehat\tau_{T,\Delta}
 \mathbin{\pm}z_{1-\eta/2}\widehat s_{T,\Delta}\bigr]\cap[-2,2],
\tag{1}
\]
with the stated quota-failure convention.  By \(\texttt{def:hybrid-interval}\), \(I_{T,\Delta}\) is instead obtained from the same logged experiment by the observable events \(E_T\) and \(G_{T,\Delta}\), using either a bounded-shift Gaussian radius containing \(B_{T,\Delta}\) or the concentration fallback.  An empirical quantile across \(200\) resampled experiments is neither construction.  Consequently, the coverage conclusions for the maps in \(\texttt{def:ordinary-wald}\) and \(\texttt{def:hybrid-interval}\) imply no coverage or tail-padding conclusion for the percentile intervals reported by Liberali et al.

Nor do the source records certify membership of the published simulations in this note's model class.  The daily Gittins-index description supplies no positive lower bound for every arm at every time.  The eta-variant supplies equal randomization only until fixed armwise minimum counts have been reached and hence supplies no permanent lower bound thereafter.  Thus the source records do not establish either of the two conditions
\[
 \pi_t(1)\geq c t^{-\alpha}
 \quad\text{and}\quad
 \pi_t(0)\geq\varepsilon
 \qquad (t\leq T),
\tag{2}
\]
which are, respectively, \(\texttt{ass:one-sided-forced-exploration}\) and \(\texttt{ass:control-floor}\).  Likewise, a reported mean time to the last observation is not the collection of armwise inequalities
\[
 P\{D_t(a)>u\}
 \leq \min\{1,L(1+u)^{-\beta}\}
 \qquad
 (a\in\mathcal A,\ u\in\mathbb N_0),
\tag{3}
\]
required by \(\texttt{ass:polynomial-delay-envelope}\).  In particular, the EUROPA mean identifies neither \(L\) nor \(\beta\), so it cannot determine the numerical bias radius in \(\texttt{def:bias-radius}\).

We next prove the limited GUSTO-1 implication.  Impose the conditions stated in the proposition: complete ascertainment of the \(30\)-day endpoint under the note's single-reveal encoding, and terminal observation for at least \(30\) days after the final randomization, so that \(\Delta_T\geq30\).  Complete ascertainment in this encoding is the support statement
\[
 P\{D_t(a)\leq30\}=1
 \qquad(a\in\mathcal A)
\tag{4}
\]
for the encoded assigned endpoint under either arm.  Since \(0<\rho<1\) and \(T\) is a positive integer, \(k_T=\lfloor T^\rho\rfloor\leq T\).  Therefore \(\texttt{def:cohort-schedule}\) gives
\[
 h_{T,\Delta}=T+\Delta_T-k_T\geq\Delta_T\geq30.
\tag{5}
\]
Equations (4)--(5) imply \(\mathbf 1\{D_t(a)>h_{T,\Delta}\}=0\) almost surely for both arms.  Substitution in \(\texttt{def:bias-radius}\) yields
\[
 \begin{aligned}
 b_{T,\Delta}
 &=P\bigl[Y(1)\mathbf 1\{D(1)>h_{T,\Delta}\}\bigr]
   -P\bigl[Y(0)\mathbf 1\{D(0)>h_{T,\Delta}\}\bigr]\\
 &=0.
 \end{aligned}
\tag{6}
\]
This conclusion uses no independence between the outcome and its delay.  It removes the delayed-tail bias only; neither (4) nor (6) supplies the allocation floors, regular-variance condition, or studentized limit theorem required for adaptive-Wald validity.  Thus complete follow-up alone does not establish coverage of (1).

Finally, consider the prospective two-arm modification described in the proposition.  Its certified independent and identically distributed potential outcome--delay vectors satisfy \(\texttt{ass:iid-arrivals}\); its binary outcomes satisfy \(\texttt{ass:bounded-outcomes}\); its predictable logged randomization satisfies \(\texttt{ass:sequential-randomization}\); its imposed permanent floors are exactly (2); and its certified tail bound is (3).  Hence, by \(\texttt{def:one-sided-policy-class}\), the resulting law and policy belong to
\(\mathfrak M_T^{\mathrm{low}}(\alpha,\beta,L,c,\varepsilon)\).  If \(Y_t(a)\) denotes the binary adverse-event indicator under arm \(a\), its causal risk difference is
\[
 P\{Y_t(1)=1\}-P\{Y_t(0)=1\}
 =P Y_t(1)-P Y_t(0)=\tau,
\tag{7}
\]
so the target is the estimand of this note rather than an odds ratio.  The uniform-honesty conclusion for \(I_{T,\Delta}\), including its unconditional quota-failure convention and concentration fallback, now follows from \(\texttt{thm:uniform-gated-inference}\).

Suppose in addition that \(\texttt{ass:regular-variance}\) holds and that the resulting regular subclass is nonempty.  By \(\texttt{prop:wald-followup-frontier}\), ordinary Wald is uniformly asymptotically honest on that subclass if and only if
\[
 x_T:=T^q(1+T+\Delta_T)^{-\beta}\longrightarrow0,
 \qquad q=(1-\alpha)/2>0.
\tag{8}
\]
The three asserted regimes follow algebraically.  If \(\beta>q\), then \(\Delta_T=0\) gives
\[
 x_T=T^{q-\beta}(1+T^{-1})^{-\beta}\longrightarrow0,
\tag{9}
\]
so no extra post-recruitment follow-up is needed.  If \(\beta=q\), positivity of \(q\) gives
\[
 x_T=\left(\frac{T}{1+T+\Delta_T}\right)^q\longrightarrow0
 \quad\Longleftrightarrow\quad
 \frac{\Delta_T}{T}\longrightarrow\infty.
\tag{10}
\]
If \(\beta<q\), positivity of \(\beta\) gives
\[
 \begin{aligned}
 x_T\longrightarrow0
 &\quad\Longleftrightarrow\quad
 \frac{1+T+\Delta_T}{T^{q/\beta}}\longrightarrow\infty\\
 &\quad\Longleftrightarrow\quad
 T+\Delta_T\gg T^{q/\beta},
 \end{aligned}
\tag{11}
\]
where the last equivalence uses \(q/\beta>1\), so \(1/T^{q/\beta}\to0\).  Equations (7)--(11) concern only the stipulated prospective two-arm risk-difference protocol after all listed conditions have been certified.  They do not transfer the result to an odds-ratio target or to the three-arm GUSTO-1 design, and (3) confirms why they yield no numerical EUROPA padding recommendation from the published mean alone.
